
\documentclass[12pt,a4paper]{article}

% ---------------------------------------------------------------
% Pacotes
% ---------------------------------------------------------------
\usepackage[utf8]{inputenc}
\usepackage[T1]{fontenc}
\usepackage[brazilian]{babel}
\usepackage[margin=2.5cm]{geometry}
\usepackage{amsmath}
\usepackage{booktabs}
\usepackage{longtable}
\usepackage{array}
\usepackage{graphicx}
\usepackage{xcolor}
\usepackage{listings}
\usepackage{enumitem}



\lstdefinestyle{c}{
  language=C,
  basicstyle=\ttfamily\small,
  keywordstyle=\color{blue!70!black},
  commentstyle=\color{green!40!black},
  stringstyle=\color{red!60!black},
  numbers=left,
  numberstyle=\tiny\color{gray},
  frame=single,
  breaklines=true,
  tabsize=4,
  showstringspaces=false,
  literate={á}{{\'a}}1 {é}{{\'e}}1 {í}{{\'i}}1 {ó}{{\'o}}1 {ú}{{\'u}}1
           {â}{{\^a}}1 {ê}{{\^e}}1 {ô}{{\^o}}1 {ã}{{\~a}}1 {õ}{{\~o}}1
           {ç}{{\c c}}1 {Á}{{\'A}}1 {É}{{\'E}}1 {Ç}{{\c C}}1
           {à}{{\`a}}1 {À}{{\`A}}1 {Í}{{\'I}}1 {Ó}{{\'O}}1 {Ú}{{\'U}}1
           {Â}{{\^A}}1 {Ê}{{\^E}}1 {Ô}{{\^O}}1 {Ã}{{\~A}}1 {Õ}{{\~O}}1
}
\lstdefinestyle{shell}{
  language={},
  basicstyle=\ttfamily\small,
  numbers=none,
  frame=single,
  breaklines=true,
  showstringspaces=false,
  literate={á}{{\'a}}1 {é}{{\'e}}1 {í}{{\'i}}1 {ó}{{\'o}}1 {ú}{{\'u}}1
           {â}{{\^a}}1 {ê}{{\^e}}1 {ô}{{\^o}}1 {ã}{{\~a}}1 {õ}{{\~o}}1
           {ç}{{\c c}}1 {Á}{{\'A}}1 {É}{{\'E}}1 {Ç}{{\c C}}1
           {à}{{\`a}}1 {À}{{\`A}}1 {Í}{{\'I}}1 {Ó}{{\'O}}1 {Ú}{{\'U}}1
           {Â}{{\^A}}1 {Ê}{{\^E}}1 {Ô}{{\^O}}1 {Ã}{{\~A}}1 {Õ}{{\~O}}1
}
\lstset{style=c}


% ---------------------------------------------------------------
% Capa
% ---------------------------------------------------------------
\title{\textbf{lib\_estat}\\[0.3em]
       \large Biblioteca Estatística em C para Variáveis Aleatórias\\[0.3em]
       \large Relatório Técnico}
\author{
  Mariana Keiko Ishiy Corrêa\\
  Maria Carolina Sobral Vila Nova de Carvalho Vieira \\
  Eddy Rodrigues Sousa \\
  Luís Guilherme de Oliveira Pelucio
}

\date{19 de outubro de 2026, Técnicas de Programação}


\begin{document}
\maketitle
\thispagestyle{empty}
\newpage

\tableofcontents
\newpage

% ===============================================================
\section{Introdução}
% ===============================================================

\subsection{Objetivo do projeto}
O projeto tem como objetivo fornecer uma biblioteca em C que implemente funções para
auxiliar o seu usuário a calcular medidas estatísticas de variáveis aleatórias, como
esperança, variância, distribuições de probabilidade e medidas de associação entre
duas variáveis aleatórias. 


\subsection{Domínio escolhido}
Variáveis aleatórias, com quatro modelos: Binomial, Geométrica, Uniforme e
Poisson, mais medidas de associação entre duas variáveis (covariância,
correlação e regressão linear).

\subsection{Escopo}
Resumo do que a biblioteca entrega: 15 funções públicas em 4 módulos, aplicação
cliente, Makefile, testes, documentação Doxygen e versão \cod{v1.0}.

\subsection{Organização do relatório}

Este relatório descreve o desenvolvimento da biblioteca \cod{lib\_estat}. Primeiro são apresentadas a divisão de tarefas, a estrutura de diretórios, a arquitetura da biblioteca, o tratamento de erros e as decisões de projeto. Em seguida, cada um dos quatro módulos é descrito, e são apresentados os testes automatizados, a depuração com GDB e a aplicação cliente. Depois vêm a compilação das bibliotecas estática e dinâmica, a documentação com Doxygen e o controle de versão. O relatório termina com as dificuldades encontradas, a conclusão e os apêndices, que reúnem o código-fonte e as saídas completas.

% ===============================================================
\section{Organização do Grupo e do Projeto}
% ===============================================================

\subsection{Divisão de tarefas}
\begin{table}[h]
\centering
\begin{tabular}{@{}p{4.2cm}p{4.2cm}p{6cm}@{}}
\toprule
Integrante & Módulo & Frente extra \\
\midrule
Mariana Keiko Ishiy Corrêa & Esperança dos modelos & Makefile, bibliotecas \cod{.a} e \cod{.so} \\

Maria Carolina Sobral Vila Nova de Carvalho Vieira & Variância dos modelos & Testes automatizados, aplicação cliente, README e montagem do relatório \\
Eddy Rodrigues Sousa & Distribuições de probabilidade & Doxygen e depuração com GDB \\
Luís Guilherme de Oliveira Pelucio & Covariância, correlação e regressão & Git (branches, merges, tag \cod{v1.0}) \\
\bottomrule
\end{tabular}
\caption{Divisão de tarefas do grupo.}
\end{table}

\subsection{Estrutura de diretórios}
\begin{lstlisting}[style=shell]
projeto2/
|-- include/   (cabecalhos .h: esperanca, variancia_modelos,
|               distribuicoes, correlacao, erros)
|-- src/       (implementacoes .c)
|-- tests/     (um test_*.c por modulo)
|-- app/       (main.c da aplicacao cliente)
|-- lib/       (libestat.a e libestat.so)
|-- build/     (objetos .o)
|-- bin/       (executaveis)
|-- docs/      (Doxygen e evidencias do GDB)
|-- Makefile
`-- README.md
\end{lstlisting}

% ===============================================================
\section{Arquitetura da Biblioteca}
% ===============================================================

\subsection{Visão geral dos módulos}

A construção do código foi feito em etapas, utilizanod cabeçalhos públicos.
\begin{lstlisting}[style=shell]
include/*.h   interfaces publicas (prototipos e comentarios Doxygen)
     |
src/*.c       implementacoes das funcoes
     |
build/*.o     objetos (build/pic/*.o para a versao dinamica)
     |
lib/          libestat.a (estatica) e libestat.so (dinamica)
     |
app/main.c, tests/*.c    programas que usam a biblioteca
     |
bin/          executaveis (app, app_dyn, test_*)
\end{lstlisting}

\subsection{Funções públicas}
\begin{table}[h]
\centering
\small
\begin{tabular}{@{}lp{9.5cm}l@{}}
\toprule
Módulo & Função & Resp. \\
\midrule
Esperança & \cod{esperanca\_binomial}, \cod{esperanca\_geometrica},
            \cod{esperanca\_uniforme}, \cod{esperanca\_poisson} & A \\
Variância & \cod{variancia\_binomial}, \cod{variancia\_geometrica},
            \cod{variancia\_uniforme}, \cod{variancia\_poisson} & B \\
Distribuições & \cod{probabilidade\_binomial}, \cod{probabilidade\_geometrica},
            \cod{probabilidade\_poisson}, \cod{densidade\_uniforme} & C \\
Correlação & \cod{covariancia}, \cod{correlacao},
            \cod{coeficiente\_linear\_regressao} & D \\
\bottomrule
\end{tabular}
\caption{As 15 funções públicas da biblioteca.}
\end{table}

\subsection{Tratamento de erros}
Todas as funções recebem um parâmetro \cod{int *erro}, preenchido com um código
definido em \cod{include/erros.h}. Em caso de erro, a função devolve um valor
neutro e informa a causa pelo código.

\begin{table}[h]
\centering
\begin{tabular}{@{}llp{7cm}@{}}
\toprule
Constante & Valor & Significado \\
\midrule
\cod{ESTAT\_OK} & 0 & Sem erro \\
\cod{ESTAT\_ERRO\_N\_INVALIDO} & 1 & Tamanho $n$ inválido (por exemplo, $n \le 0$) \\
\cod{ESTAT\_ERRO\_TAMANHOS\_DIFERENTES} & 2 & Vetores de tamanhos diferentes (reservado: nenhuma função o utiliza) \\
\cod{ESTAT\_ERRO\_DIVISAO\_ZERO} & 3 & Divisão por zero (por exemplo, $p = 0$ na geométrica) \\
\cod{ESTAT\_ERRO\_PARAMETRO\_INVALIDO} & 4 & Parâmetro fora do domínio (por exemplo, $p \notin [0,1]$, $a \ge b$ ou vetor \cod{NULL}) \\
\bottomrule
\end{tabular}
\caption{Códigos de erro da biblioteca.}
\end{table}


O ponteiro \cod{erro} não pode ser \cod{NULL}, com uma exceção: as funções do módulo de correlação aceitam \cod{erro == NULL}, caso em que nenhum código de erro é armazenado. 
Nesse módulo, um vetor \cod{NULL} resulta em \cod{ESTAT\_ERRO\_PARAMETRO\_INVALIDO}.

\subsection{Decisões de projeto}
\begin{itemize}[leftmargin=*]
  \item Erro via parâmetro \cod{int *erro} (e não via \cod{errno} ou \cod{exit}),
        para que a biblioteca nunca encerre o programa do usuário, uma vez que a biblioteca é uma ferramenta de apoio, e não um programa completo.

  \item Cabeçalhos independentes, cada um incluindo apenas \cod{erros.h}.
  \item Convenção da distribuição geométrica: $k = 0, 1, 2, \ldots$ (número de falhas até o primeiro sucesso).
  Uma vez que há uma variação em cada representação da distribuição geométrica, decidimos adotar a convenção de $k$ como o número de falhas até o primeiro sucesso, que é a forma mais intuitiva, no estudo de variáveis aleatórias.
 
  \item Tratamento de $p = 0$ na geométrica: tanto \cod{esperanca\_geometrica} quanto \cod{variancia\_geometrica} devolvem \cod{ESTAT\_ERRO\_DIVISAO\_ZERO}, pois as fórmulas $(1-p)/p$ e $(1-p)/p^{2}$ dividem por $p$. Valores de $p$ fora de $[0,1]$ resultam em \cod{ESTAT\_ERRO\_PARAMETRO\_INVALIDO}.
\end{itemize}

% ===============================================================
\section{Descrição dos Módulos}
% ===============================================================

% ---------------------------------------------------------------
\subsection{Esperança dos modelos}
% ---------------------------------------------------------------
\label{sec:esperanca}
Este módulo calcula a esperança dos quatro modelos (geométrico, uniforme, binomial e Poisson). Os protótipos estão em
\cod{include/esperanca.h} e a implementação em \cod{src/esperanca.c}.

\subsubsection{Fórmulas}
\begin{table}[h]
\centering
\begin{tabular}{@{}lcc@{}}
\toprule
Modelo & Parâmetros & Esperança \\
\midrule
Binomial   & $n, p$    & $E[X] = n\,p$ \\
Geométrica & $p$       & $E[X] = \dfrac{1-p}{p}$ \\
Uniforme   & $a, b$    & $E[X] = \dfrac{a+b}{2}$ \\
Poisson    & $\lambda$ & $E[X] = \lambda$ \\
\bottomrule
\end{tabular}
\caption{Esperança dos modelos implementados.}
\end{table}

\subsubsection{Assinaturas das funções utilizadas}
\begin{lstlisting}
double esperanca_binomial(int n, double p, int *erro);
double esperanca_geometrica(double p, int *erro);
double esperanca_uniforme(double a, double b, int *erro);
double esperanca_poisson(double lambda, int *erro);
\end{lstlisting}

\subsubsection{Validações}
\begin{table}[h]
\centering
\small
\begin{tabular}{@{}lll@{}}
\toprule
Função & Condição & Código de erro \\
\midrule
\cod{esperanca\_binomial}   & $n \le 0$          & \cod{ESTAT\_ERRO\_N\_INVALIDO} \\
                            & $p < 0$ ou $p > 1$ & \cod{ESTAT\_ERRO\_PARAMETRO\_INVALIDO} \\
\cod{esperanca\_geometrica} & $p < 0$ ou $p > 1$ & \cod{ESTAT\_ERRO\_PARAMETRO\_INVALIDO} \\
                            & $p = 0$            & \cod{ESTAT\_ERRO\_DIVISAO\_ZERO} \\
\cod{esperanca\_uniforme}   & $a \ge b$          & \cod{ESTAT\_ERRO\_PARAMETRO\_INVALIDO} \\
\cod{esperanca\_poisson}    & $\lambda \le 0$    & \cod{ESTAT\_ERRO\_PARAMETRO\_INVALIDO} \\
\bottomrule
\end{tabular}
\caption{Validações de entrada do módulo de esperança.}
\end{table}

\subsubsection{Exemplos numéricos}
\begin{table}[h]
\centering
\begin{tabular}{@{}lc@{}}
\toprule
Chamada & Resultado \\
\midrule
\cod{esperanca\_binomial(10, 0.5)} & 5 \\
\cod{esperanca\_geometrica(0.5)}   & 1 \\
\cod{esperanca\_uniforme(5, 10)}   & 7,5 \\
\cod{esperanca\_poisson(10)}       & 10 \\
\bottomrule
\end{tabular}
\caption{Valores conferidos pelos testes e pela aplicação cliente.}
\end{table}
% ---------------------------------------------------------------
\subsection{Variância dos modelos}
\label{sec:variancia}
% ---------------------------------------------------------------
Este módulo calcula a variância dos quatro modelos de variáveis aleatórias. Os protótipos estão em
\cod{include/variancia\_modelos.h} e a implementação em
\cod{src/variancia\_modelos.c}.

\subsubsection{Fórmulas}
\begin{table}[h]
\centering
\begin{tabular}{@{}lcc@{}}
\toprule
Modelo & Parâmetros & Variância \\
\midrule
Binomial   & $n, p$        & $\operatorname{Var}(X) = n\,p\,(1-p)$ \\
Geométrica & $p$           & $\operatorname{Var}(X) = \dfrac{1-p}{p^{2}}$ \\
Uniforme   & $a, b$        & $\operatorname{Var}(X) = \dfrac{(b-a)^{2}}{12}$ \\
Poisson    & $\lambda$     & $\operatorname{Var}(X) = \lambda$ \\
\bottomrule
\end{tabular}
\caption{Variância dos modelos implementados.}
\end{table}

Na geométrica foi usada a fórmula fechada $(1-p)/p^{2}$, e não o cálculo a partir
de $E[X^2] - E[X]^2$, porque a primeira é mais eficiente e não há risco de subtração de números próximos que causem perda de precisão.


\subsubsection{Assinaturas das funções utilizadas}
\begin{lstlisting}
double variancia_binomial(int n, double p, int *erro);
double variancia_geometrica(double p, int *erro);
double variancia_uniforme(double a, double b, int *erro);
double variancia_poisson(double lambda, int *erro);
\end{lstlisting}

\subsubsection{Validações}
\begin{table}[h]
\centering
\small
\begin{tabular}{@{}lll@{}}
\toprule
Função & Condição & Código de erro \\
\midrule
\cod{variancia\_binomial}   & $n \le 0$             & \cod{ESTAT\_ERRO\_N\_INVALIDO} \\
                            & $p < 0$ ou $p > 1$    & \cod{ESTAT\_ERRO\_PARAMETRO\_INVALIDO} \\
\cod{variancia\_geometrica} & $p < 0$ ou $p > 1$    & \cod{ESTAT\_ERRO\_PARAMETRO\_INVALIDO} \\
                            & $p = 0$               & \cod{ESTAT\_ERRO\_DIVISAO\_ZERO} \\
\cod{variancia\_uniforme}   & $a \ge b$             & \cod{ESTAT\_ERRO\_PARAMETRO\_INVALIDO} \\
\cod{variancia\_poisson}    & $\lambda \le 0$       & \cod{ESTAT\_ERRO\_PARAMETRO\_INVALIDO} \\
\bottomrule
\end{tabular}
\caption{Validações de entrada do módulo de variância.}
\end{table}

\subsubsection{Exemplo de implementação}
\begin{lstlisting}
double variancia_binomial(int n, double p, int *erro){
    if (n <= 0) {
        *erro = ESTAT_ERRO_N_INVALIDO;
        return 0.0;
    }
    if (p < 0 || p > 1) {
        *erro = ESTAT_ERRO_PARAMETRO_INVALIDO;
        return 0.0;
    }

    *erro = ESTAT_OK;
    return n*p*(1-p);


}
\end{lstlisting}

\subsubsection{Exemplos numéricos}
\begin{table}[h]
\centering
\begin{tabular}{@{}lc@{}}
\toprule
Chamada & Resultado \\
\midrule
\cod{variancia\_binomial(10, 0.5)}   & 2,5 \\
\cod{variancia\_geometrica(0.5)}     & 2 \\
\cod{variancia\_uniforme(0, 10)}     & 8,3333 \\
\cod{variancia\_poisson(4)}          & 4 \\
\bottomrule
\end{tabular}
\caption{Valores conferidos manualmente e pela aplicação cliente.}
\end{table}


% ---------------------------------------------------------------
\subsection{Distribuições de probabilidade}
% ---------------------------------------------------------------
\todo{autor: C. Fórmulas, validações, convenção de $k$ na geométrica, exemplos}

% ---------------------------------------------------------------
\subsection{Covariância, correlação e regressão}
% ---------------------------------------------------------------
\todo{autor: D. Fórmulas, vetores de tamanhos diferentes, desvio padrão zero, exemplos}

% ===============================================================
\section{Testes Automatizados}
% ===============================================================

\subsection{Padrão adotado}
Os testes usam \cod{assert} da biblioteca padrão, com uma função auxiliar de
comparação com tolerância para números em ponto flutuante:
\begin{lstlisting}
int aproximadamente(double x, double y) {
    return fabs(y - x) < 0.000001;
} 
\end{lstlisting}
Escolheu-se \cod{assert} por ser simples, confiável e amplamente usado, e não criar macros próprias que poderiam confundir o leitor.
Também utilizou-se a função \cod{aproximadamente} para comparar os números double e evitar problemas de arredondamento que podem ocorrer 
ao longo do código.


Cada módulo tem um arquivo \cod{tests/test\_<modulo>.c}. Quando um teste passa, o programa imprime uma mensagem no terminal (no módulo de correlação, uma por grupo de testes). 
Quando um teste falha, o \cod{assert} aborta o programa e o \cod{make} interrompe a execução.


\subsection{Estratégia de casos}
Cada função é testada com (1) caso normal com valor conhecido, (2) cada
condição de erro, e (3) valores limite ($p = 0$, $p = 1$, $n = 0$ etc.).

\subsection{Cobertura}
\begin{table}[h]
\centering
\small
\begin{tabular}{@{}lcc@{}}
\toprule
Módulo & Funções & Testes \\
\midrule
Esperança      & 4 & 10 \\
Variância      & 4 & 15 \\
Distribuições  & 4 & 25 \\
Correlação     & 3 & 16 \\
\midrule
Total          & 15 & 71 \\
\bottomrule
\end{tabular}
\caption{Cobertura dos testes por módulo.}
\end{table}

Nesta tabela, cada função de teste dos arquivos \cod{tests/test\_*.c} conta como um teste.

\subsubsection{Testes do módulo de variância}
\begin{table}[h]
\centering
\small
\begin{tabular}{@{}llp{6cm}@{}}
\toprule
Função & Nº de testes & Casos cobertos \\
\midrule
Binomial   & 6 & normal, $n$ inválido, $n = 0$, $p$ inválido, $p = 0$, $p = 1$ \\
Geométrica & 4 & normal, $p$ inválido, divisão por zero, $p = 1$ \\
Uniforme   & 3 & normal, $a \ge b$, $a$ negativo \\
Poisson    & 2 & normal, $\lambda$ inválido \\
\bottomrule
\end{tabular}
\caption{Testes de \cod{variancia\_modelos}.}
\end{table}

\subsection{Como executar}
\begin{lstlisting}[style=shell]
make test
\end{lstlisting}
A saída abaixo foi obtida com todos os módulos passando (arquivo \cod{saida\_test.txt}, gerado por \cod{make -s test}):
\saida{saida_test.txt}

% ===============================================================
\section{Depuração com GDB}
% ===============================================================

Cenário: introduzir um defeito proposital, observar o teste falhar, localizar a
causa com o depurador e corrigir.

\subsection{Estado inicial (testes passando)}
\todo{saída de \cod{evidencia\_1\_passa.txt}}

\subsection{Defeito introduzido}
Troca do sinal em \cod{variancia\_binomial}:
\begin{lstlisting}
return n * p * (1 + p);   /* defeito: o correto e (1 - p) */
\end{lstlisting}
\todo{diff de \cod{evidencia\_2\_bug.diff}}

\subsection{Falha do teste}
Para $n = 10$ e $p = 0{,}5$ o valor esperado é $2{,}5$, mas o código devolve
$7{,}5$.
\todo{saída de \cod{evidencia\_3\_falha.txt} com função e linha do \cod{assert}}

\subsection{Sessão de depuração}
\begin{table}[h]
\centering
\small
\begin{tabular}{@{}lp{7.5cm}@{}}
\toprule
Comando & Objetivo \\
\midrule

\cod{gdb ./bin/test\_variancia\_modelos} & abrir o executável compilado com \cod{-g} \\
\cod{break variancia\_binomial}         & parar na função suspeita \\
\cod{run}                               & executar até o ponto de parada \\
\cod{info args}                         & ver $n$ e $p$ recebidos \\
\cod{next}                              & avançar uma linha \\
\cod{print n*p*(1-p)}                   & valor da fórmula correta \\
\cod{print n*p*(1+p)}                   & valor da fórmula do código \\
\cod{bt}                                & pilha de chamadas \\
\bottomrule
\end{tabular}
\caption{Comandos usados na depuração.}
\end{table}
\todo{transcrição de \cod{evidencia\_4\_depuracao.txt} (sessão no Linux)}

\subsection{Correção e verificação}
\todo{diff de correção e saída com os testes voltando a passar}

\subsection{Lições}
\todo{2 ou 3 frases sobre o que o depurador mostrou que a mensagem do
\cod{assert} sozinha não mostrava}

% ===============================================================
\section{Aplicação Cliente}
% ===============================================================

\subsection{Objetivo}
Mostrar o uso da biblioteca por um programa externo, que inclui apenas os
cabeçalhos públicos e nunca os \cod{.c}.

\subsection{Estrutura do programa}
Funções \cod{demo\_esperanca}, \cod{demo\_variancia}, \cod{demo\_distribuicoes}
e \cod{demo\_correlacao}, mais as auxiliares \cod{mensagem\_erro} e
\cod{imprime\_res}, que traduzem o código de erro em texto.

\subsection{Chamadas inválidas de propósito}
Cada demonstração inclui chamadas com argumentos inválidos para mostrar o
tratamento de erros (por exemplo, \cod{variancia\_binomial(10, 1.5)}).

\subsection{Saída da aplicação}
A execução de \cod{./bin/app} produz a saída abaixo (arquivo \cod{saida\_app.txt}):
\saida{saida_app.txt}

\subsection{Versões estática e dinâmica}
O uso e a compilação de bibliotecas estáticas e dinâmicas são demonstrados no \cod{Makefile}. A aplicação cliente é compilada duas vezes: uma com a biblioteca estática (\cod{libestat.a}) e outra com a biblioteca dinâmica (\cod{libestat.so}).

No executável compilado pela biblioteca estática, a biblioteca é incorporada ao executável durante a compilação da aplicação, assim permitindo que a biblioteca seja utilizada mesmo que não esteja presente no sistema do usuário.

No executável compilado pela biblioteca dinâmica, o código não é incorporado: o executável guarda apenas a referência à biblioteca, e o carregador dinâmico a procura quando o programa é executado. 
Se ele não a encontrar, ocorre um erro de carregamento da biblioteca. Para evitar isso, a pasta onde ela está pode ser indicada pela variável de ambiente \cod{LD\_LIBRARY\_PATH} ou gravada no próprio executável
na hora da ligação, com a opção \cod{-Wl,-rpath}. 
Já as opções \cod{-L} e \cod{-l} são usadas apenas na compilação, para que o linker localize a biblioteca.


Os executáveis gerados são \cod{bin/app} (estático) e \cod{bin/app\_dyn} (dinâmico), e a solução adotada pelo grupo é a variável \cod{LD\_LIBRARY\_PATH=lib}. A parte dinâmica foi verificada em Linux, e não no macOS. Sem a variável, o \cod{ldd} não encontra a biblioteca e a execução falha. Com a variável, o \cod{ldd} a localiza em \cod{lib/} (trecho relevante da saída):
\begin{lstlisting}[style=shell]
$ ldd bin/app_dyn
        libestat.so => not found
        libm.so.6 => /lib64/libm.so.6
        libc.so.6 => /lib64/libc.so.6

$ LD_LIBRARY_PATH=lib ldd bin/app_dyn
        libestat.so => lib/libestat.so
        libm.so.6 => /lib64/libm.so.6
        libc.so.6 => /lib64/libc.so.6

$ ./bin/app_dyn
./bin/app_dyn: error while loading shared libraries: libestat.so: cannot open shared object file: No such file or directory
\end{lstlisting}


% ===============================================================
\section{Compilação e Distribuição}
% ===============================================================

\subsection{Makefile}
O Makefile usa variáveis para os diretórios (\cod{SRC\_DIR}, \cod{BUILD\_DIR},
\cod{LIB\_DIR}, \cod{BIN\_DIR}) e regras genéricas com \cod{wildcard} e
\cod{\%}, tal que um novo arquivo em \cod{src/} ou \cod{tests/} é
incluído sem precisar alterar o Makefile.

\begin{table}[h]
\centering
\small
\begin{tabular}{@{}ll@{}}
\toprule
Alvo & Função \\
\midrule
\cod{make}          & gera \cod{lib/libestat.a} e \cod{lib/libestat.so} \\
\cod{make test}     & compila e executa os testes com a biblioteca estática \\
\cod{make test-dyn} & compila e executa os testes com a biblioteca dinâmica \\
\cod{make app}      & gera \cod{bin/app} (estático) \\
\cod{make app-dyn}  & gera \cod{bin/app\_dyn} (dinâmico) \\
\cod{make clean}    & remove os arquivos gerados \\
\cod{make run}      & compila e executa \cod{bin/app} (estático) \\
\cod{make run-dyn}  & compila e executa \cod{bin/app\_dyn} com \cod{LD\_LIBRARY\_PATH} \\
\bottomrule
\end{tabular}
\caption{Alvos do Makefile.}
\end{table}

\subsection{Biblioteca estática e dinâmica}
\begin{itemize}[leftmargin=*]
  \item Estática: os \cod{.c} viram objetos (\cod{.o}) e são empacotados com
        \cod{ar rcs lib/libestat.a build/*.o}. O código usado é copiado para o
        executável, que não depende da biblioteca na execução.
  \item Dinâmica: os objetos são compilados com \cod{-fPIC} (código
        independente de posição) e unidos com \cod{gcc -shared}. O executável
        guarda só a referência, e a biblioteca é carregada na execução.
  \item Ambas são ligadas com \cod{-lm}, colocado no final do comando por causa
        da ordem de resolução dos símbolos.
\end{itemize}

\subsection{Compilação incremental}
Como os objetos dependem do seu \cod{.c} e dos cabeçalhos, alterar um
\cod{.c} recompila apenas ele e refaz as bibliotecas.
\begin{lstlisting}[style=shell]
gcc -Wall -Wextra -Iinclude -c src/esperanca.c -o build/esperanca.o
ar rcs lib/libestat.a build/correlacao.o build/distribuicoes.o ...
gcc -Wall -Wextra -Iinclude -fPIC -c src/esperanca.c -o build/pic/esperanca.o
gcc -shared build/pic/correlacao.o ... -o lib/libestat.so -lm
\end{lstlisting}

\subsection{Localização da biblioteca dinâmica}
Na ligação, \cod{-L} indica o diretório da biblioteca e \cod{-l} o seu nome
(\cod{-lestat} procura \cod{libestat.so}). Porém, na execução o carregador
não procura em \cod{lib/}, e o programa falha com erro de biblioteca não
encontrada. A solução adotada é a variável de ambiente:
\begin{lstlisting}[style=shell]
LD_LIBRARY_PATH=lib ./bin/app_dyn
\end{lstlisting}
Alternativamente, o caminho pode ser gravado no executável com
\cod{-Wl,-rpath}.

% ===============================================================
\section{Documentação com Doxygen}
% ===============================================================

\subsection{Convenções de comentário}
Todos os cabeçalhos usam \cod{@file}, \cod{@brief}, \cod{@param} (incluindo
\cod{erro}) e \cod{@return}.

\subsection{Geração}
\begin{lstlisting}[style=shell]
doxygen Doxyfile
\end{lstlisting}
\todo{autor: C. Mostrar a página inicial e uma página de função no HTML}

\subsection{Verificação}
A geração sobre o diretório \cod{include/} foi executada sem nenhum aviso de documentação
(\cod{WARN\_IF\_UNDOCUMENTED}, \cod{WARN\_NO\_PARAMDOC} ativos). Durante essa
verificação, foi detectada e corrigida uma declaração duplicada de
\cod{variancia\_uniforme} no cabeçalho do módulo de variância.

% ===============================================================
\section{Controle de Versão}
% ===============================================================

\subsection{Fluxo de branches}
\todo{autor: D. Branches \cod{main}, \cod{develop}, \cod{esperanca}, \cod{modelos},
\cod{variancia}, \cod{app\_client} etc., e a regra de integração em \cod{develop}}

\subsection{Histórico}
\todo{\cod{git log --oneline --graph} resumido}

\subsection{Conflitos e como foram resolvidos}
\todo{se houve algum}

\subsection{Versão \cod{v1.0}}
\todo{tag, data e hash do commit}

% ===============================================================
\section{Dificuldades e Soluções}
% ===============================================================
\begin{table}[h]
\centering
\small
\begin{tabular}{@{}p{5cm}p{9cm}@{}}
\toprule
Dificuldade & Solução \\
\midrule
Convenção da geométrica diferente entre módulos & 
A esperança usava $1/p$ (número de tentativas), divergindo dos demais módulos (número de falhas), e $p = 0$ gerava códigos de erro diferentes. O grupo adotou $k$ como número de falhas ($E[X]=(1-p)/p$) e \cod{ESTAT\_ERRO\_DIVISAO\_ZERO} para $p = 0$ em ambos os módulos. A correção foi aplicada à esperança (commit \cod{1ba33a9}). \\
Declaração duplicada no cabeçalho de variância & Detectada pelo Doxygen e removida \\
GDB indisponível no macOS com chip Apple & 
Sessão feita em máquina Linux; ensaio prévio com \cod{lldb} \\
Nome da biblioteca gerada (\cod{liblib\_estat}), com o prefixo \cod{lib} duplicado & Corrigido no Makefile (\cod{LIB\_NAME = estat}), que passou a gerar \cod{libestat.a} e \cod{libestat.so} \\
Biblioteca dinâmica não encontrada na execução & Uso de \cod{LD\_LIBRARY\_PATH=lib} (alternativa: \cod{-Wl,-rpath}) \\
\bottomrule
\end{tabular}
\caption{Principais problemas encontrados.}
\end{table}

% ===============================================================
\section{Conclusão}
% ===============================================================
O projeto entregou uma biblioteca em C com 15 funções públicas distribuídas em quatro módulos (esperança, variância, distribuições de probabilidade e correlação), com tratamento de erros padronizado pelo parâmetro \cod{int *erro}. 
Também foram entregues as versões estática (\cod{libestat.a}) e dinâmica (\cod{libestat.so}), geradas por um Makefile, os testes automatizados de todos os módulos, uma aplicação cliente que demonstra o uso das funções e a documentação gerada com Doxygen.

Durante o trabalho, o grupo aprendeu a padronizar convenções entre módulos (como a da distribuição geométrica e o código de erro para $p = 0$), a comparar números em ponto flutuante com tolerância, a diferença entre ligação estática e dinâmica e a usar branches para integrar o trabalho de várias pessoas. 
Como melhorias futuras, a biblioteca poderia incluir novos modelos, como a Normal e a Exponencial, e a parte dinâmica poderia ser validada também no macOS.

% ===============================================================
\appendix
\section{Código-fonte}
% ===============================================================
\subsection{\cod{include/variancia\_modelos.h}}
\fonte{../../include/variancia_modelos.h}

\subsection{\cod{src/variancia\_modelos.c}}
\fonte{../../src/variancia_modelos.c}

\subsection{\cod{tests/test\_variancia\_modelos.c}}
\fonte{../../tests/test_variancia_modelos.c}

\subsection{\cod{app/main.c}}
\fonte{../../app/main.c}

\section{Saídas completas}

As saídas completas do \cod{make test} e da aplicação cliente estão nas seções \emph{Testes Automatizados} e \emph{Aplicação Cliente}.

\section{README}

O repositório contém um \cod{README.md} na raiz, com a introdução, as instruções de compilação, execução e testes, a convenção de erros e um exemplo de uso. 
Ele também é usado como página principal da documentação gerada pelo Doxygen.

\end{document}